% Ad Familiares 5.2 --- headnote
% ad-familiares-05-02, January or early February 62 BC, written at Rome

\textit{Cicero's reply to Q. Caecilius Metellus Celer (Fam
5.1), written at Rome at the very end of January or the
beginning of February 62 BC. The longest piece of self-
defence in the early letter corpus, and the closest thing
we have to Cicero's own narrative of the Metellus Nepos
clash. Cicero takes Celer's complaint sentence by sentence
and answers it with patient, even grave courtesy. The
joke in the senate that Celer's relations had got him to
suppress the praise he meant to give Cicero --- the joke
which had set Celer off --- is acknowledged: it was made,
the laughter was at Cicero's own expectation more than at
Celer, the underlying expectation was honourable. The
``mutual feeling'' Celer invoked is examined: ``what you
reckon mutual in friendship I do not know; I reckon it
this: when an equal goodwill is received and returned.''}

\textit{The heart of the letter (\textsection 6--9) is
the narrative of the Nepos affair. Nepos as tribune, on
the day before the Kalends of January, had stripped
Cicero of the power to address the people on the last
day of his consulship; Cicero in reply swore in a great
voice the truest and finest of oaths, which the people
swore back at him in turn that he had truly sworn ---
the famous oath ``that the commonwealth and this city
had been saved by me alone'' (Pis. 6, off. 3.74). On the
Kalends of January and three days later in the senate
Cicero withstood Nepos directly: ``to his rashness, if I
had not stood up by virtue and by spirit, who would not
have judged that I had been brave in my consulship by
chance rather than by counsel?'' The closing paragraph
(\textsection 10) is the model of a Roman peace-keeping
in a quarrel: Cicero pardons Celer's grief, even praises
it (``my own feeling reminds me how great is the force
of brotherly love''), and asks Celer to be a fair judge
of his own. The breach was not healed, and Celer would
ally with Clodius against Cicero in the years to come.}
